\section{Discussion}
\label{discussion}

The results and explanations raise three questions: why word counts hold their own against a
pretrained transformer, what the one-class SVM tells us about fraud in EMSCAD, and how far these
conclusions reach. This section takes them in turn.

\subsection{Why word counts compete with a pretrained transformer}

TF-IDF logistic regression matches fine-tuned DistilBERT, and in the stack it carries most of the
weight. Our interpretation is that much of the fraud signal in EMSCAD lies in exact words: pay
amounts, the templates of a few posters, employer names and the placeholders that replaced
contact details. A sparse linear model gives each such word its own weight and can learn it from a
handful of the 514 training frauds. DistilBERT, by contrast, reads only the first 256 tokens and
represents an ad by its meaning, which helps where the wording is new (the 12 frauds that only
DistilBERT catches) but adds little where a known phrase already decides. On this dataset, then,
the transformer is most useful as a second opinion that makes different mistakes, which is exactly
what the stack exploits; on its own it costs a GPU without improving on the linear model.

\subsection{What the one-class SVM shows}

The one-class SVM began as a near-failure and became an informative model once its input and
kernel width were set properly, and both of its behaviours say something about the data. With a
broad kernel it ranked frauds as \emph{more} typical than real ads. Many frauds are repeated
templates, so they sit in dense regions of the text space, and an anomaly detector that judges ads
against the bulk of the data treats them as normal. With a local kernel it asks a different
question, whether an ad has close neighbours among the real ads, and that question separates the
classes reasonably well (ROC-AUC 0.86), in our reading because a fraud template rarely has a real
near-copy.

Even at its best, however, it reaches a PR-AUC of only 0.176, against 0.66 for a supervised model
on the same features, so labels remain the scarce resource that matters: a job board without
labelled fraud would catch some frauds with anomaly detection, but far from all. Its real value in
this study is that it fails differently. It catches 16 seed-0 test frauds that no supervised model
catches, mostly hotel and call-centre vacancies that every supervised model scored below its
threshold. Its low precision makes it unsuitable on its own, but as an additional signal for a
moderator, or as a third input to the stack, it covers a blind spot of the supervised models.

\subsection{Limitations}

Several limits qualify these conclusions. All results come from a single dataset, and the country
artefact and the poster templates show that part of what the models learn reflects how EMSCAD was
collected, so the numbers would not transfer unchanged to another job board. The tuning budget of
the supervised models is small (4 random candidates with 3-fold cross-validation per seed), a
deliberate choice to keep the classical runs to about an hour on a laptop; a larger search could
move individual models, but every supervised model received the same budget.

The validation rows also serve more than one purpose. DistilBERT chooses its best epoch on them and
then its threshold on the same rows, so its validation scores are slightly optimistic, and the
stacker is fitted on those same scores. The test rows, however, are never used for any choice, so
the test numbers remain honest. For the one-class SVM the training score is meaningless, because
its local kernel memorises the real training ads.

Finally, the five splits reuse the same 17,599 ads, so the standard deviations measure sensitivity
to the split rather than sampling error over new ads. The explanations are local: Integrated
Gradients covers ten ads from one seed and the coefficient table one refit, so they show what
drives these cases rather than a general rule. DistilBERT also sees only the first 256 tokens of
each ad, whereas the TF-IDF models see every word.
